\section{Conclusion}
\label{sec:conclusion}

This study evaluated continual personalization for Mandarin VSR under a strict predict-then-update protocol. Across three seeds, no-feedback adaptation of a 75k-parameter residual adapter reduced mean Chinese-LiPS CER by 1.59 absolute points relative to the frozen source model. Supplying one correction every ten utterances added a further 0.832-point mean gain, with the same direction in every matched seed, and produced a 2.42-point total gain over static. These are two valid results under different supervision budgets and should not be collapsed into one claim.

The experiments do not support the current structural implementation. The original dynamic expert bank improved over one adapter by only 0.055 CER points on average and assigned 99.87\% of test-revisit samples to one expert. A validation-only calibration passed its signature proxy gate but then created eight experts, reused the A1-dominant expert on only 42.3\% of A2, and performed significantly worse than one adapter. Because all expert runs could use the known feedback-location flag for pre-prediction growth, they are retained only as pre-fix diagnostics and do not show that dynamic experts are intrinsically ineffective. Reliability and rollback ablations likewise yielded effects too small for a strong causal safety claim.

Validation rejects three superficially plausible update rules. Greedy correct-span preservation is significantly worse than whole-sequence replay, while adapter-only entropy minimization degrades CER by more than 24 points. Target-conditioned error localization significantly outperforms same-mass randomized support, confirming that correction placement carries information, but remains 2.367 CER points worse than replay and also worsens revisit difference-in-differences. It therefore does not enter test evaluation.

On the disjoint target-dev3 development stream, pseudo-only and feedback-only adaptation each significantly improve over static, but adding both update sources does not significantly improve over either source alone. Under identical 75-query budgets, uncertainty selection is 0.272 CER points better than periodic but misses the 0.3-point threshold and has a confidence interval crossing zero; it is 0.021 points worse than random on the point estimate and does not enrich the binary true-error rate. Its static-corrected forgetting interval also fails the pre-registered non-degradation condition. This active-query route is therefore a development no-go, as is the target-dev2 hybrid route. Holdout2 remains frozen and unread, and neither failed route receives extra seeds or tuning.

Before submission, the study still requires stronger matched standard TTA and VSR-personalization baselines, a second real shift type beyond speakers, and a new pre-registered candidate that passes development before any confirmatory holdout run. The resolved gates remove the current expert-growth, error-localized, hybrid, and uncertainty-query implementations from the proposed deployment path. The supported system remains the simpler sparse-feedback single adapter, with its measured resource footprint and limitations; failed candidates are retained as diagnostic evidence rather than promoted through repeated tuning on consumed streams.
